\subsection{Peto's Paradox and Cancer Resistance}
Peto's paradox highlights the intriguing observation that cancer risk does not scale with body size across species, despite the increased number of cell divisions in larger organisms \citep{peto1977epidemiology}. This paradox has led to the discovery of various cancer resistance mechanisms in large-bodied animals, including elephants with their multiple TP53 gene copies \citep{sulak2016tp53} and naked mole-rats with enhanced extracellular matrix components \citep{tian2013high}. However, the evolutionary strategies employed by other large mammals, particularly domesticated species like pigs, remain less explored.

\subsection{The TP53 C-terminal Domain as an Evolutionary Hotspot}
The TP53 tumor suppressor protein is a critical regulator of cell cycle arrest, DNA repair, and apoptosis in response to cellular stress. While the DNA-binding domain is highly conserved across vertebrates, the C-terminal domain (CTD) exhibits remarkable evolutionary plasticity. The CTD serves as a regulatory hub, containing multiple post-translational modification (PTM) sites that modulate TP53's stability, DNA-binding affinity, and protein-protein interactions \citep{kruse2009mcp}. This domain's variability across species suggests it may serve as an evolutionary "tuning knob" for TP53 function.

\subsection{Hypothesis and Study Aims}
We hypothesize that pigs (\textit{Sus scrofa}) have evolved a unique "less-is-more" strategy in their TP53 CTD, characterized by selective degeneration of regulatory elements while maintaining core tumor suppressor function. To test this hypothesis, we conducted a comprehensive comparative analysis of TP53 sequences across eleven mammalian species, with particular focus on:

\begin{enumerate}
    \item Quantifying CTD-specific GC content and selection pressure (dN/dS ratios)
    \item Mapping PTM site conservation and loss patterns
    \item Analyzing structural consequences of CTD degeneration
    \item Correlating CTD evolutionary patterns with species-specific cancer resistance
\end{enumerate}

Our findings reveal that pigs have undergone significant CTD degeneration while maintaining functional TP53 activity, representing a novel evolutionary solution to Peto's paradox with potential implications for understanding cancer resistance mechanisms and TP53 regulation in humans.
